
\section{Preesquema con grupo finito de operadores; preesquema cociente}
\label{V.1}

Sean $X$ un preesquema y $G$ un grupo finito que actúa sobre~$X$ por
automorfismos, por la derecha para fijar ideas. Si $X$ es afín de anillo $A$,
$G$ actúa, por tanto, por automorfismos a la izquierda sobre~$A$.

Para todo preesquema $Z$, $G$ actúa a la izquierda sobre el conjunto
$\Hom(X,Z)$; por consiguiente, puede considerarse el conjunto
$$
\Hom(X,Z)^G
$$
de los morfismos invariantes por~$G$. Este conjunto depende funtorialmente
de~$Z$, y cabe preguntarse si este funtor es «representable», 
\ie isomorfo a un funtor $Z\mto\Hom(Y,Z)$.
\marginpar{106}

Esto significa que pueden hallarse un preesquema $Y$ y un morfismo invariante
por~$G$
$$
p\colon X\to Y
$$
tales que, para todo~$Z$, la aplicación correspondiente $g\mto gp$
$$
\Hom(Y,Z)\to \Hom(X,Z)^G
$$
sea biyectiva. Se dice entonces que $(Y,p)$ es un
\emph{preesquema cociente}
\index{cociente (preesquema)|hyperpage}%
de~$X$ \ifthenelse{\boolean{orig}}{por~$g$}{por~$G$}
(está determinado salvo isomorfismo único).

\begin{proposition}
\label{V.1.1}
Sean $A$ un anillo sobre el que el grupo finito $G$ actúa a la izquierda,
$B=A^G$ el subanillo de invariantes de~$A$, $X=\Spec(A)$ e $Y=\Spec(B)$, y
sea $p\colon X\to Y$ el morfismo canónico (evidentemente invariante
por~$G$). Entonces:
\begin{enumerate}
\item[(i)] $A$ es entero sobre~$B$, \ie $p$ es un morfismo
\emph{entero}.
\item[(ii)] El morfismo $p$ es sobreyectivo, sus fibras son las órbitas
de~$G$ y la topología de~$Y$ es el cociente de la de~$X$.
\item[(iii)] Sean $x\in X$, $y=p(x)$ y $G_x$ el estabilizador de~$x$.
Entonces
\ifthenelse{\boolean{orig}}
{$\kappa(x)$ es una extensión algebraica cuasigaloisiana de $\kappa(y)$ y
la aplicación canónica de~$G_x$ en el grupo
$\Gal(\kappa(x)/\kappa(y))$ de los $\kappa(y)$-automorfismos de
$\kappa(x)$}
{$\kres(x)$ es una extensión algebraica cuasigaloisiana de $\kres(y)$ y
la aplicación canónica de~$G_x$ en el grupo
$\Gal(\kres(x)/\kres(y))$ de los $\kres(y)$-automorfismos de
$\kres(x)$}
es \ifthenelse{\boolean{orig}}{sobreyectivo.}{sobreyectiva.}
\item[(iv)] $(Y,p)$ es un preesquema cociente de~$X$ por~$G$.
\end{enumerate}
\end{proposition}

Los enunciados
\ifthenelse{\boolean{orig}}
{(i) (ii) (iii)}
{(i), (ii), (iii)}
son bien conocidos en álgebra conmutativa\footnote{\Cf N. Bourbaki,
Alg. Comm. Chap\ptbl~5, \S~1 y \S~2, th\ptbl~2.} y solo se recuerdan aquí,
salvo la afirmación sobre la topología, que procede del siguiente hecho
general, consecuencia fácil del teorema de Cohen--Seidenberg: un morfismo
entero es cerrado (\ie transforma conjuntos cerrados en conjuntos cerrados).
Observemos enseguida lo siguiente:

\begin{corollaire}
\label{V.1.2}
Bajo las condiciones anteriores, el homomorfismo natural
$$\cal{O}_Y{\to} p_*(\cal{O}_X)^G$$
es un isomorfismo.
\end{corollaire}

Esto resulta inmediatamente de la fórmula
$$
(S^{-1}A)^G = S^{-1}(A^G),
$$
válida para toda parte multiplicativamente estable~$S$ de~$B=A^G$
(fórmula que admite una versión con módulos y que se enuncia, más
generalmente, para un cambio de base $A\to A'$ que sea \emph{plano}),
\ifthenelse{\boolean{orig}}{aplicado}{aplicada}
al caso en que~$S$ está engendrada por un elemento~$f$ de~$B$.

La afirmación~(ii) y el cor\ptbl~\Ref{V.1.2} implican fácilmente~(iv); más
generalmente, se tiene lo siguiente:

\begin{proposition}
\label{V.1.3}
Sean
\marginpar{107}
$X$ un preesquema con un grupo finito de automorfismos~$G$ y
$p\colon{}X\to{}Y$ un morfismo afín invariante tal que
$\cal{O}_Y\isomto p_*(\cal{O}_X)^G$. Entonces siguen siendo válidas las
conclusiones
\ifthenelse{\boolean{orig}}
{\textup{(i) (ii) (iii) (iv)}}
{\textup{(i), (ii), (iii), (iv)}}
de~\Ref{V.1.1}.
\end{proposition}

En efecto, para
\ifthenelse{\boolean{orig}}
{(i) (ii) (iii)}
{(i), (ii), (iii),}
puede suponerse que~$Y$, y por tanto~$X$, es afín; si $B,A$ son sus anillos,
la hipótesis implica $B=A^G$, y basta aplicar~\Ref{V.1.1}. Para~(iv), se
utilizan~(ii) y $\cal{O}_Y=p_*(\cal{O}_X)^G$.

\begin{corollaire}
\label{V.1.4}
Bajo las condiciones de~\Ref{V.1.3}, para todo abierto~$U$ de~$Y$, $U$ es
un cociente de $X|U=p^{-1}(U)$ por~$G$.
\end{corollaire}

En efecto, $p^{-1}(U)\to{}U$, inducido por~$p$, satisface las mismas
hipótesis que~$p$.

Si ahora~$X$ es un $Z$-preesquema y las operaciones de~$G$ son
$Z$-automorfismos, entonces, por~(iv), $Y$ es un $Z$-preesquema. Dicho esto:

\begin{corollaire}
\label{V.1.5}
Para que $X$ sea afín, \resp separado sobre~$Z$, es necesario y suficiente
que~$Y$ lo sea. Si $X$ es de tipo finito sobre~$Z$, es \emph{finito}
sobre~$Y$; si además~$Z$ es localmente noetheriano, $Y$ es de tipo finito
sobre~$Z$.
\end{corollaire}

Como $X$ es afín y, con mayor razón, separado sobre~$Y$, si $Y$ es afín,
\resp separado sobre~$Z$, también lo es~$X$. Recíprocamente, supongamos que
$X$ es afín sobre~$Z$ y probemos que~$Y$ lo es: gracias a~\Ref{V.1.4} puede
suponerse que~$Z$ es afín, y la cuestión se reduce a demostrar que si~$X$ es
afín, también lo es~$Y$, lo cual resulta de la determinación explícita de~$Y$
como
\ifthenelse{\boolean{orig}}{$\Spec(A)^G$}{$\Spec(A^G)$}
efectuada en~\Ref{V.1.1}.
Del mismo modo, como $p\colon{}X\to{}Y$ es entero, y por tanto
universalmente cerrado, y es sobreyectivo, se sigue que, si~$X$ es separado
sobre~$Z$, también lo es~$Y$ (¡lema que convendría enunciar por separado!);
en efecto, en el diagrama
$$
\xymatrix@C=1.5cm{ X \times_ZX\ar[r]^-{p\times_Zp} & Y\times_ZY \\
X \ar[u]^{\Delta_{X/Z}}\ar[r]^{p} & Y \ar[u]_{\Delta_{Y/Z}} }
$$
el morfismo $X\times_ZX\to Y\times_ZY$ es cerrado, por lo que transforma la
diagonal (cerrada) de~$X\times_ZX$ en una parte cerrada de~$Y\times_ZY$, que
no es sino la diagonal de este último, pues~$p$ es sobreyectivo.
--- Si $X$ es de tipo finito sobre~$Z$, lo es con mayor razón sobre~$Y$ y,
por consiguiente, es finito sobre~$Y$ (pues el morfismo $X\to Y$ ya es
entero).
Supongamos además que~$Z$ es localmente noetheriano y probemos que~$Y$ es de
tipo finito sobre~$Z$. Gracias a~\Ref{V.1.4} puede suponerse que~$Z$ es afín.
Como el espacio topológico~$X$ es cuasicompacto y
$p\colon X\to Y$ es sobreyectivo, $Y$ es
\ifthenelse{\boolean{orig}}{tambien}{también}
cuasicompacto y, por tanto, unión finita de abiertos
\marginpar{108}
afines; por~\Ref{V.1.4}, la cuestión se reduce al caso en que~$Y$ es afín y,
por tanto, también~$X$. Ahora bien, el anillo~$A$ de~$X$ es un álgebra de
tipo finito sobre el anillo noetheriano~$C$ de~$Z$, y se sabe que
$B=A^G$ también es un álgebra de tipo finito sobre~$C$ (pues $A$ será
entera, y por tanto finita, sobre una subálgebra~$B'$ de~$B$ de tipo finito
sobre~$C$; como~$B'$ es noetheriana, el álgebra~$B$ también es
\ifthenelse{\boolean{orig}}{finito}{finita}
sobre~$B'$, y por consiguiente es de tipo finito sobre~$C$).
